\section{The auxiliary coordinate is a uniformizer in the unchanged original DVR}
\begin{lemma}\label{lem:original-whole-trait-open-lift}
A morphism from the spectrum of a local ring factors through every actual open immersion whose image contains the image of its closed point. The resulting entire trait has exactly the original projection.
\lean{EnvelopeProperty.SchemeOriginalAdaptedTrait.lift, EnvelopeProperty.SchemeOriginalAdaptedTrait.lift_projection}
\leanok
\end{lemma}
\begin{proof}
Every prime of a local ring specializes to its closed point. Its image under the original morphism therefore specializes to the image of that closed point. An open set containing the latter contains all its generalizations. Thus the entire original trait lies in the image of the open immersion. The universal property of that actual open immersion supplies the unique factorization and its exact projection equality.
\end{proof}

\begin{theorem}\label{thm:normalized-actual-original-uniformizer-trait}
Under the original characteristic-zero normalized coefficient-slice hypotheses, let the original parameter $t$ be nonzero and belong to the maximal ideal of the original DVR $V$. There is an actual affine adapted etale chart of the unchanged normalized model containing its original center. The entire original trait factors through this chart, with exactly its original projection. Its canonical native center-stalk map to the unchanged $V$ is bijective. Its actual first coordinate evaluates to an irreducible element $\pi\in V$, and the whole auxiliary parameter diagram is the polynomial evaluation map at this same $\pi$. There are a unit $v\in V$ and an unrestricted positive integer $e$ with $t=v\pi^e$ and original native order $\operatorname{ord}_V(t)=e$.
\uses{thm:normalized-actual-smooth-auxiliary-parameter,lem:original-whole-trait-open-lift,thm:normalized-slice-original-center-ring,thm:point-closure-original-stalk-maximal-ideal}
\lean{EnvelopeProperty.SchemeOriginalAdaptedTrait.exists_original_uniformizer_trait}
\leanok
\end{theorem}
\begin{proof}
Take the actual original adapted etale chart and factor the entire original trait through its open immersion by the preceding lemma. Its closed point still maps to the unchanged original center. The actual original residual ideal at that center is the entire native maximal ideal. An open immersion induces a native stalk isomorphism, so its pulled-back residual ideal at the lifted center is again the entire native maximal ideal. Adaptedness identifies this ideal with the ideal generated by the germ of the actual first coordinate.

The original canonical center-stalk map to $V$ is bijective. Its composite description through the open-immersion stalk map is exactly the canonical local map from the lifted trait. The open-immersion stalk map is itself bijective, so bijectivity of this composite proves bijectivity of the lifted canonical local map. A surjective local map takes the entire maximal ideal onto the entire maximal ideal. Apply it to the preceding principal ideal equality. The actual value $\pi$ of the first coordinate thus generates the actual maximal ideal of the unchanged original DVR. It is consequently irreducible.

The native center-stalk triangle identifies the whole lifted map to the coordinate plane with the ring homomorphism obtained by the same canonical local map on every coordinate. The original coefficient square makes this homomorphism preserve every element of $E$. Composing with the actual first-coordinate polynomial map therefore gives evaluation at the same element $\pi$, on every original polynomial, as asserted.

Finally write the unchanged original nonzero $t$ in the DVR as $v\pi^e$ with $v$ a unit. If $e=0$, this makes $t$ a unit, contradicting its original maximal-ideal membership. Hence $e>0$. A unit has order zero and this irreducible element has order one, so the original native order of $t$ is exactly $e$. This conclusion allows every positive multiplicity and changes neither $V$ nor its original fraction field.
\end{proof}
